\section{ Parameters for $PSL(2,\C)$-characters of surfaces}
\label{PARAM} Fix a compact closed oriented surface $S$ of genus
$g$ with a non empty set $V= \{v_1,\dots \ , v_r\}$ of marked
points, and negative Euler characteristic $\chi(S\setminus V)<0$.
Denote by $\pi$ the fundamental group of $S\setminus V$, and by
$$\Rr(g,r) = {\rm Hom}(\pi, PSL(2,\C))/PSL(2,\C)$$ the set of {\it
all} conjugacy classes of $PSL(2,\C)$-valued representations of
$\pi$. The group $\pi$ is free of rank $\kappa = 2g + r -1$. Any
choice of free generators of $\pi$ identifies the set ${\rm
Hom}(\pi, PSL(2,\C))$ with $PSL(2,\C)^\kappa$. Different such
identifications are related by algebraic automorphisms of
$PSL(2,\C)^\kappa$.  Moreover, the isomorphism $PSL(2,\C)\cong
SO(3,\C)$ induced by the adjoint action ${\rm Ad}: PSL(2,\mc)
\rightarrow {\rm Aut}(sl(2,\mc))$ implies that ${\rm Hom}(\pi,
PSL(2,\C))$ is an {\it affine} complex algebraic set, with the
complex algebraic action of $PSL(2,\C)$ by conjugation.

As the rough topological quotient space $\Rr(g,r)$ is not
Hausdorff, it is more convenient to consider the algebro-geometric
quotient $$X(\pi)={\rm
  Hom}(\pi, PSL(2,\C))//PSL(2,\C)$$ of invariant theory, called the
{\it variety of $PSL(2,\C)$-characters} (see eg. \cite{HP}).
Recall that $X(\pi)$ is a complex affine algebraic set such that
there exists a surjective regular map $$t:{\rm Hom}(\pi,
PSL(2,\C))\to X(\pi)$$ inducing an isomorphism $t^*$ between the
regular functions on $X(\pi)$ and the regular functions on ${\rm
Hom}(\pi, PSL(2,\C))$ invariant by conjugation. In general
$t(\gamma)=t(\sigma)$ does not imply that $\gamma$ and $\sigma$
are conjugate, but this is true if we restrict to {\it
irreducible} representations: we have ${\rm
  Hom}^{irr}(\pi, PSL(2,\C)) = t^{-1}(X^{irr}(\pi))$, where
$X^{irr}(\pi)= t({\rm Hom}^{irr}(\pi, PSL(2,\C))$, so that the
(restricted) rough quotient $\Rr(g,r)^{irr}$ and the algebraic
quotient $X^{irr}(\pi)$ coincide.

 We can deal with the whole of $\Rr(g,r)$ to construct
the QHFT. Hence below we treat the complex dimension of subsets of
$\Rr(g,r)$ somewhat formally, as everything can be substantiated
in terms of $X(\pi)$, or by restriction to $X^{irr}(\pi)$. For
instance, as $PSL(2,\C)$ has trivial centre and complex dimension
$3$ we say that the complex dimension of $\Rr(g,r)$ is $3\kappa -
3 = -3\chi(F)$.
\subsection {Efficient triangulations}\label{e-triang}
Fix a surface $F$ with $r$ boundary components, obtained by
removing from $S$ the interior of small $2$-disks $D_i$ such that
$v_i \in \partial D_i$.

A triangulation $T'$ of $S$ with the set of vertices equal to $V$ is
called a {\it topological ideal triangulation} of $S\setminus
V$. Given such a $T'$, we need a marking of corners of the
$2$-simplices. The best suited to $3$-dimensional extension are
induced by global branchings $b'$ of $T'$, and it is known that pairs
$(T',b')$ always exist. Then, as in section \ref{QHG}, we have a sign
function $\sigma = \sigma_{(T',b')}$. A {\it corner map} $v\mapsto
c_v$ associates to each vertex $v$ of $T'$ the corner at $v$ of a
triangle, say $t_v$, in its star. We say that $v\mapsto c_v$ is {\it
$t$-injective} if $v\mapsto t_v$ is injective.
\begin{lem}\label{injcorner}
For every $(g,r)\ne (0,3)$, every triangulation $T'$ of $S$ with $r$
vertices admits $t$-injective corner maps.
\end{lem}
\noindent {\it Proof.} For $(g,r) =(0,4)$ or $g>0$ and $r=1$, it
is immediate to construct such a triangulation. Subdividing a
triangle by taking the cone from an interior point preserves the
existence of $t$-injective corner maps. By induction on $r$ we
deduce that for every $(g,r)\ne (0,3)$ there exist triangulations
of $S$ as in the statement.\hfill\break In Figure \ref{eTflip} the
corner selection is specified by a $*$, and the rows show all the
possible flip moves on triangulations of $S$, up to obvious
symmetries, that preserve the injectivity of corner maps. Consider
triangulations $T'$, $T''$ of $S$ with $r$ vertices, such that
$T''$ supports a $t$-injective corner map. It is well known that
$T''$ is connected to $T'$ via a finite sequence of flips. The
$t$-injective corner map for $T''$ yields a $t$-injective corner
map for $T'$ by decorating these flips as in Figure
\ref{eTflip}.\hfill$\Box$\medskip

\begin{figure}[ht]
\begin{center}
\includegraphics[width=6cm]{eTflip.eps}
\caption{\label{eTflip} The flips with marked corners.}
\end{center}
\end{figure}
There are no obstructions to use arbitrary corners maps in what
follows. Specializing to injective ones just simplifies the
exposition. For $(g,r)= (0,3)$, triangulations have two
$2$-simplices, one with two selected corners.  From now on, we
assume that $(g,r) \neq (0,3)$, the extension to the $(0,3)$ case
being straightforward.

\smallskip

Given a pair $(g,r) \ne (0,3)$ and $(T',b')$ as in Lemma
\ref{injcorner}, fix a $t$-injective corner map $v\mapsto c_v$. In the
interior of each triangle $t_v$ consider a bigon $D_v$ with one vertex
at $v$, and call $v'$ the other vertex. Remove from $t_v$ the interior
of $D_v$, and triangulate the resulting cell $s_v = t_v \setminus {\rm
  Int}(D_v)$ by taking the cone with base $v'$. Repeating this
procedure for each $t_v$, we get a triangulation $T$ of $F$ with $2r$
vertices and $p+2r$ triangles, where $p$ denotes the number of
triangles of $T'$.  The set of edges of $T$, $E(T)$, contains $E(T')$
in a natural way, and $|E(T)|= |E(T')|+4r$. We extend $b'$ to a
branching $b$ on $T$ as in in Figure \ref{newb}, and the sign function
$\sigma_{(T',b')}$ to the triangles of $(T,b)$ in the natural way.
Note that the figure shows only one of the possible branching
configuration.  In general we extend $b'$ to $b$ so that we can
recover $(T',b')$ from $(T,b)$ by ``zipping'' and ``collapsing'', as
suggested at the bottom of Figure \ref{newb}. In what follows, for
simplicity we will refer to this configuration, as the treatment of
the others is similar.
\begin{figure}[ht]
\begin{center}
\includegraphics[width=6cm]{newb2.eps}
\caption{\label{newb} The branched triangulated cell $s_v$.}
\end{center}
\end{figure}
\begin{defi}\label{defet}
{\rm We call the pair $(T,b)$ an efficient triangulation (for short:
    $e$-{\it triangulation}) of $F$. For each vertex $v$, the {\it
    preferred innner triangle at} $v$ is the triangle $\tau_v$ in
    $s_v$ with an edge on $\partial F$ whose $b$-orientation coincides
    with the boundary orientation of $F$.}
\end{defi}

\subsection{Cocycle parameters}\label{Dpar}
The inclusions of Int$(F)$ into $F$ and $S\setminus V$ induce
identifications of the respective fundamental groups.

Fix an $e$-triangulation $(T,b)$ of $F$. Denote by $Z(T,b)$ the space
of $PSL(2,\C)$-valued $1$-cocycles on $(T,b)$; we stipulate
that on a triangle with ordered $b$-oriented edges $e_0,e_1,e_2$ the
cocycle relation is $z(e_0)z(e_1)z(e_2)^{-1} = 1$. In this section we
construct a parametrization of $\Rr(g,r)$ based on cocycle
coefficients by specifying subsets of $Z(T,b)$ with small `residual
gauge groups', that make principal algebraic bundles over the
``strata'' of a suitable partition of $\Rr(g,r)$. These strata are
determined by the holonomies around the boundary components of
$F$.

Write $C(T,b)$ for the space of $PSL(2,\C)$-valued $0$-cochains on
$(T,b)$, that is the $PSL(2,\C)$-valued functions defined on the
set of vertices of $T$.  Two $1$-cocycles $z$ and $z'$ are said
equivalent {\it up to gauge transformation} if there is a
$0$-cochain $\lambda$ such that, for every oriented edge
$e=[x_0,x_1]$, we have $$z'(e) = \lambda(x_0)^{-1}z(e)\lambda(x_1)
\ . $$ It is well known that the quotient set $H(T,b) =
Z(T,b)/C(T,b)$ is in one-one correspondence with $\Rr(g,r)$.
Indeed, for any fixed $x_0 \in T$ we have a natural surjective map
$$f_{x_0}: Z(T,b) \to {\rm Hom}(\pi, PSL(2,\C)),$$ where
$\pi=\pi_1(F,x_0)$, and two representations $f_{x_0}(z)$ and
$f_{x_0}(z')$ define the same point in $\Rr(g,r)$ if and only if
$z$ and $z'$ are equivalent up to gauge transformation. Note that
the complex dimension of $H(T,b)$ is $3(|E(T)| - (p+2r) - 2r) =
-3\chi(F)$, which is the dimension of $\Rr(g,r)$. This is because
the complex dimension of $C(T,b)$ is equal to $6r$, the set
$Z(T,b)$ is defined by $3(p+2r)$ polynomial relations on $3|E(T)|$
variables, and $PSL(2,\C)$ has trivial centre. (Recall that $p$ is
the number of triangles of the initial triangulation $T'$ of $S$
with $r$ vertices.)

Denote by $B^+(2,\C)$ (respectively $B^-(2,\C)$) the Borel subgroup of
$SL(2,\C)$ of upper (respectively lower) triangular matrices. Let
$PB^\pm(2,\C)= B^\pm(2,\C)/\pm I$ and put
$$ P=\left( \begin{array}{cc} 0 & 1 \\ 1 & 0  \end{array}\right).$$
Define a map
\begin{equation}\label{convert}
\Psi : SL(2,\C)\to SL(2,\C),\quad \Psi(A)=PAP
\end{equation}
and denote by $\Psi$ the induced automorphism of $PSL(2,\C)$. We
have $\Psi(B^\pm(2,\C)) = B^\mp(2,\C)$. %We note that for cocycle
%parameters, the roles of $B^\pm(2,\C)$ are perfectly exchangeable
%each with the other via the map $\Phi$. This symmetry will be
%broken in Subsection \ref{Ipar}, when we will consider the
%$\Ii$-parameters.
For any $g\in PSL(2,\C)$ we distinguish the {\it type} of $g$ as
{\it
  trivial}, {\it parabolic} or {\it generic} (the latter for elliptic or
loxodromic) with the obvious meaning. We denote
$$\mathcal{C}^+(g) \in PB^+(2,\C)$$ the canonical
upper triangular matrix (up to sign) representative of the conjugacy
class of $g$ (for generic $g$ we normalize $\mathcal{C}^+(g)$ by stipulating
that the top diagonal entry has absolute value $>1$), and we set
$$\mathcal{C}^-(g) = \Psi(\mathcal{C}^+(g)). $$
We define the type of $\rho \in \Rr(g,r)$ as the $n$-uple of types of
the $\rho$-holonomies of the oriented boundary components
$\gamma_1,\ldots,\gamma_r$ of $F$. Put:
$$\begin{array}{l} \Rr(g,r,t) = \{ \rho \in \Rr(g,r) \mid \rho {\rm\
has \ type\ } t\} \\ \Rr(g,r,\mathcal{C}^\pm) = \{ \rho \in \Rr(g,r,t)
\mid \mathcal{C}^\pm(\rho(\gamma_i)) = \mathcal{C}_i^\pm\}
\end{array}$$
where $\mathcal{C}^\pm=(\mathcal{C}_1^\pm,\ldots,\mathcal{C}_r^\pm)$,
$\mathcal{C}_i^\pm$ is an arbitrary diagonal or unipotent element in
$PB^\pm(2,\C)$, and $\mathcal{C}^\pm$ has type $t$.
\smallskip
For any $g \in PB^\pm(2,\C)$ let us write $g=[a,b]^\pm$, where $a$ is
the top diagonal entry of $g$ and $b$ is the non diagonal one (we do
the abuse of confusing $g$ with its projective class). For every
vertex $v_i$ of $(T,b)$, let $e_1^i$ be the boundary edge of the
preferred inner triangle $\tau_{v_i}$ (see Definition \ref{defet}),
and $\gamma_i$ the oriented boundary loop of $F$ based at
$v_i$. Recall the projections $f_{v_i}: Z(T,b) \to {\rm Hom}(\pi,
PSL(2,\C))$.
\begin{defi} \label{defcoc}
The set of ($\pm$)-{\it cocycle parameters} for
$\Rr(g,r,\mathcal{C}^\pm)$ is $$\Zz(T,b,\mathcal{C}^\pm) = \{ z
\in Z(T,b) \mid \forall i=1,\ldots,r,\ f_{v_i}(z)(\gamma_i) =
\mathcal{C}^\pm_i,\ z(e_1^i)= [1,1/2]^\pm \}.$$
\end{defi}
Clearly $\Zz(T,b,\mathcal{C}^\pm)$ is non empty. In fact, any
$z\in Z(T,b)$ with ${\rm conj} \circ f(z) \in
\Rr(g,r,\mathcal{C}^\pm)$ is equivalent to one in
$\Zz(T,b,\mathcal{C}^\pm)$ via some gauge transformation.
Moreover, given a vertex $v_i$ and a $0$-cochain $s$ with support
at $v_i$ and $v_i'$, $s$ maps $\Zz(T,b,\mathcal{C}^\pm)$ onto
itself if and only if $$s(v_i)\in {\rm Stab}(\mathcal{C}^\pm_i)$$
and $$s(v_i')^{-1}[1,1/2]^\pm s(v_i) = [1,1/2]^\pm. $$ As the
triangles $\tau_{v_i}$ are in one-one correspondence with the
$v_i$, we deduce that there is a projection
\begin{equation}\label{projc}
p^\pm_{\mathcal{C}}: \Zz(T,b,\mathcal{C}^\pm)\to
\Rr(g,r,\mathcal{C}^\pm) \end{equation} with fiber isomorphic to
the group $$\Gg (T,b,\mathcal{C}^\pm):={\rm
Stab}(\mathcal{C}^\pm_1)\times \dots \times {\rm
Stab}(\mathcal{C}^\pm_r)$$ that we call the {\it residual gauge
transformations}.  Denote by $\Zz(T,b,t)^\pm$ (respectively
$p^\pm_t$) the union of the $\Zz(T,b,\mathcal{C}^\pm)$
(respectively $p^\pm _{\mathcal{C}}$) over all $\mathcal{C}^\pm$
with type $t$, and put $$\textstyle \Zz(T,b)^\pm = \coprod_t
\Zz(T,b,t)^\pm, \quad p^\pm =  \coprod_t p^\pm_t).$$ For futur
reference, we summarize the above constructions in the following
proposition.
\begin{prop}\label{topcomp} For every $\Cc^\pm$, the projection
  $p^\pm : \Zz(T,b)^\pm \rightarrow \Rr(g,r)$ restricts
to a complex affine algebraic principal $\Gg
(T,b,\mathcal{C}^\pm)$-bundle $$\Zz(T,b,\mathcal{C}^\pm) \to
\Rr(g,r,\mathcal{C}).$$ The map $\Psi$ in (\ref{convert}) yields
an isomorphism $\Zz(T,b,\mathcal{C}^\pm) \cong
\Zz(T,b,\mathcal{C}^\mp)$. Moreover we have: $$\begin{array}{lll}
{\rm dim}(\Zz(T,b,\mathcal{C}^\pm)= -3\chi(F)\\ {\rm
dim}(\Rr(g,r,\mathcal{C})) = -3\chi(F) - {\rm dim}(\Gg
(T,b,t_{\mathcal{C}}))\\ {\rm dim}(\Rr(g,r,t_{\mathcal{C}})) =
{\rm dim}(\Rr(g,r,\mathcal{C})) + \alpha(t_{\mathcal{C}})
\end{array}$$
where $t_{\mathcal{C}}$ is the type of $\mathcal{C}$ and
$\alpha(t_{\mathcal{C}})$ is the number of
generic entries of $t_{\mathcal{C}}$.
\end{prop}
Observe that: if $t_{{\rm gen}}$ is the purely generic type, then
${\rm dim}(\Rr(g,r,t_{{\rm gen}}))= -3\chi(F)$; if $t_{{\rm par}}$ is
purely parabolic, then ${\rm dim}(\Rr(g,r,t_{{\rm par}}))= -3\chi(F)-
r$; if $t_I$ is purely trivial, then ${\rm dim}(\Rr(g,r,t_I))=
-3\chi(F) -3r = -3\chi(S)= 6g-6$. Hence $\Rr(g,r,t_{{\rm gen}})$ is a
dense open subset of $\Rr(g,r)$. Moreover, via the inclusion of
closures, we have a filtration of $\Rr(g,r)$ for which $\Rr(g,r,t_I)$
is the `deepest' part. It would be interesting to study the
singularities of the closure of each $\Rr(g,r,t)$ in $\Rr(g,r)$, in
order to check if this filtration induces a ``stratification'' of
$\Rr(g,r)$.

\begin{remark}\label{real} {\rm For parabolic elements $g \in
    PSL(2,\R)$ we have two conjugacy classes, that can be
distinguished by a sign. So, replacing $PSL(2,\C)$ with
$PSL(2,\R)$, the constructions of this section still work by
associating a {\it sign} to each parabolic end of the surface $S
\setminus V$. Note that for $PSL(2,\R)$-valued cocycles the
idealization procedure described in Section \ref{Ipar} below gives
only degenerate triangles and tetrahedra (with real shapes or
cross-ratio moduli).}
\end{remark}

\noindent {\bf Example: the Fricke space.} Suppose that $S$ has a
unique marked point $v$. Choose a standard curve system
$\mathcal{S}=\{a_i,b_i\}_{i=1}^{g}$ based at $v$, so that
$$\pi_1(S,v) =<a_1,b_1,\ldots,a_g,b_g\mid [a_1,b_1]\ldots
[a_g,b_g] =1>.$$ Cutting open $S$ along $\mathcal{S}$ we get a
$4g$-gon $P$ with oriented boundary edges. Taking the cone to a
vertex it is easy to construct a branched triangulation of $P$,
which induces one, say $(T',b')$, for $(S,v)$.  Denote by $(T,b)$
any $e$-triangulation obtained from $(T',b')$. \hfill\break Recall
that the Teichm\"uller space $\mathcal{T}(S)$ can be identified
with the set of conjugacy classes of $PSL(2,\mr)$-valued discrete
faithful representations of $\pi_1(S)$, and that we have the
well-known (real-analytic) Fricke parametrization $\mathcal{T}(S)
\cong \mr^{6g-6}$ (see e.g. \cite{Ab}). For each $z \in
\mathcal{T}(S)$ the Fricke coordinates of $z$ are matrix entries
of the $\hat{z}(\gamma)$ for all $\gamma \in \mathcal{S}$, where
$\hat{z}$ is a representative of $z$ specified by fixing once and
for all three of the fixed points of $z(a_g)$ and $z(b_g)$. So
$\mathcal{T}(S)$ embeds in the space of real cocycle parameters
for $F$ with trivial type $t_I$. This embedding is generalized
easily to the Teichm\"uller space of arbitrary bordered Riemann
surfaces, by considering the spaces $\Zz(T,b,t)$ for all types $t$
and $e$-triangulations of $(S,V)$ with arbitrary $V$ (see Remark
\ref{real}).

\subsection {Cross-ratio parameters}\label{Ipar}
In this section we derive from the cocycle parameters
$\Zz(T,b)^\pm$ other parameters for $\Rr(g,r)$ which are related
to the shear-bend coordinates for pleated hyperbolic surfaces.
These parameters are obtained via an {\it idealization procedure}
that includes the choice of a base point on the Riemann sphere. We
fix this base point as $0$. Note that the $B^+(2,\C)$-orbit of $0$
is the whole of $\C$, while $B^-(2,\C)$ fixes $0$. Hence the
symmetry between $\Zz(T,b)^+$ and $\Zz(T,b)^-$ given by the map
$\Psi$ in (\ref{convert}) shall be broken.
\begin{defi} \label{triideal-pre} {\rm
(Compare with Definition \ref{3dim-id}.) Let $(K,b)$ be any oriented
surface branched triangulation. Let $z$ be any $PSL(2,\C)$-valued
cocycle on $(K,b)$.  We say that $z$ is {\it idealizable} if for any
triangle $t$ of $(T,b)$ with $b$-ordered edges $e_0$, $e_1$ and $e_2$,
the points $u_0=0$, $u_1= z(e_0)(0)$ and $u_2= z(e_2)(0)$ are distinct
in $\mc$. We say that the complex triangle with vertices $u_0$, $u_1$
and $u_2$ is the \emph{idealization} of $t$.}
\end{defi}
\begin{defi} \label{triideal}
{\rm Let $(T,b)$ be an $e$-triangulation of $F$ obtained from a
branched ideal triangulation $(T',b')$ of $S \setminus V$. A
cocycle $z\in \Zz(T,b)^+$ is {\it (strongly) idealizable} if:
\smallskip

(a) $z$ is idealizable;
\smallskip

(b) $\Psi(z) \in \Zz(T,b)^-$, which {\it cannot} be idealizable at the
triangles having an edge on $\partial F$, is nevertheless idealizable
at every other triangle.
\smallskip

We denote by $\Zz_I(T,b)^+$ the set of (strongly) idealizable
cocycles, and we put $\Zz_I(T,b)^- = \Psi(\Zz_I(T,b)^+)$ and
$\Rr_I(T,b) = p^+(\Zz_I(T,b)^+) = p^-(\Zz_I(T,b)^- )$.}
\end{defi}

Clearly, $\Zz_I(T,b)^+$ is a non-empty dense open subset of
$\Zz(T,b)^+$. If every edge of $T$ has distinct endpoints (in case
we say that $T$ {\it quasi-regular}), then $\Rr_I(T,b) =
\Rr(g,r)$. In general, characters of representations with a free
action on a non-empty domain of $\mc\mathbb{P}^1$ (such as
quasi-Fuchsian representations) always belong to $\Rr_I(T,b)$, for
any $(T,b)$. By using the arguments of \cite{Kmod}, Theorem 1, it
can be shown that for any character $\rho\in \Rr(g,r)$ of
irreducible representations, there exists an $e$-triangulation
$(T,b)$ with idealizable cocycles representing $\rho$. In fact,
the union of a finite number of spaces $\Zz_I(T,b)^+$ cover the
whole of $\Rr(g,r)^{irr}$.

\smallskip

\noindent{\bf Exponential $\Ii$-parameters.}  For any cocycle $z
\in \Zz_I(T,b)^+$, we associate a non zero complex weight
$W^+(z)(e)$ to each edge $e$ of $T$ that is not contained in
$\partial F$, as follows. Let $p_e$ be the initial endpoint of
$e$, and $t_l$ and $t_r$ the left and right adjacent triangles (as
viewed from $e$). Locally modify the branching on $t_l \cup t_r$
by cyclically reordering the vertices on each triangle, so that
$p_e$ is eventually the source of the new branching on both $t_l$
and $t_r$. The {\it ($+$)-exponential $\Ii$-parameter} $W^+(z)(e)$
is the cross-ratio modulus at $e$ of the (possibly degenerate)
branched oriented hyperbolic ideal tetrahedron spanned by the
idealization of $t_r \cup t_l$, where the branching completes the
one of $t_l \cup t_r$ so that $*_b=1$. %With the terminology of
\smallskip

Let us assume now that $\Psi(z)\in \Zz_I(T,b)^-$. For each $v \in V$,
denote by $e_v^1$ and $e_v^2$ the edges of the triangle $t_v$ of $T'$
having $c_v$ as corner. At an edge $e$ of $T'$ distinct from any of
the $e_v^i$, we define the {\it ($-$)-exponential $\Ii$-parameter}
$W^-(\Psi(z))(e)$ in the same way as $W^+(z)(e)$, but taking
the idealization for $\Psi(z)$ instead of $z$. If $e$ is one of the
$e_v^i$, the formula works as well, except that each left/right
triangle with an edge on $\partial F$ is replaced with the innermost
triangle in the corresponding triangle $t_v$ of $T'$. Note that $t_r
\cup t_l$ is again a quadrilateral, because $0$ is a fixed point of
the cocycle values at the boundary edges of $F$.
\smallskip

By varying the edge in $T$ or $T'$ for every $\Cc^\pm$, we have two
{\it ($\pm$)-exponential $\Ii$ parameter maps}
(recall that the number of edges of $T'$ is $ -3\chi(F)$, and $p$
denotes the number of triangles of $T'$):
$$W^+: \Zz_I(T,b,\mathcal{C}^+)
    \rightarrow (\mc \setminus \{ 0\})^{-3\chi(F)+2p}$$
$$W^-: \Zz_I(T,b,\mathcal{C}^-) \rightarrow (\mc \setminus \{ 0
    \})^{-3\chi(F)}.$$

    \begin{defi}\label{phasespace} {\rm We call
        $W^\pm(T,b,\mathcal{C}) = W^\pm(\Zz_I(T,b,\mathcal{C}^\pm))$
        the {\it ($\pm$)-exponential $\Ii$-parameter space} of
        $\Rr_I(T,b,\mathcal{C})$. }
\end{defi}
We stress again that the result of the idealization strongly
depends on the choice of the base point, here $0$. In particular,
the exponential $\Ii$-parameters are {\it not} invariant under
arbitrary gauge transformations of cocycles.  A remarkable
exception is for gauge transformations associated to $0$-cochains
$\lambda$ with values in $PB^-(2,\C)$. Indeed, these act on the
idealization of quadrilaterals as conformal transformations of the
{\it four} vertices (this is because $0$ is fixed by every
$\lambda(v)$, $v \in V$), and cross-ratios are conformal
invariants. This makes a big difference between $W^+$ and $W^-$.
In fact the whole of the residual gauge group
$\Gg(T,b,\mathcal{C}^+)$ acts on $W^+(T,b,\mathcal{C})$ via the
map $W^+$. On the other hand, consider the subgroup of
$\Gg(T,b,\mathcal{C}^-)$ defined as $$B\Gg(T,b,\mathcal{C}^-):=
\prod_{i=1,\dots, r}{\rm Stab}(\mathcal{C}^-_i)\cap PB^-(2,\C) \
.$$ Any fiber $(W^-)^{-1}(W^-(z))$ is given by the
$B\Gg(T,b,\mathcal{C}^-)$-orbit of $z$ in
$\Zz_I(T,b,\mathcal{C}^-)$. Hence the actual residual gauge
transformations of $W^-(T,b,\mathcal{C})$ are in one-one
correspondence with the quotient set $\Gg(T,b,\mathcal{C}^-)/
B\Gg(T,b,\mathcal{C}^-)$, where the equivalence class of $\lambda$
is $B\Gg(T,b,\mathcal{C}^-)\lambda$. For the matter of notational
convenience, formally put $\textstyle B\Gg(T,b,\mathcal{C}^+):=
\prod_{i=1,\dots, r} {\rm Id}$. The map $$\Theta^\pm:
W^\pm(T,b,\mathcal{C}) \to \Rr(g,r,\mathcal{C}^\pm)$$ given by $
\Theta^\pm(W^\pm(z))= p^\pm_{\mathcal{C}}(z)$, with
$p^\pm_{\mathcal{C}}$ defined in (\ref{projc}), is a principal
$\Gg (T,b,\mathcal{C}^\pm)/ B\Gg(T,b,\mathcal{C}^\pm)$-bundle. The
situation is particularly clean when $\mathcal{C}$ has no trivial
entries :
\begin{prop}\label{no-gauge} If $\mathcal{C}$ has no trivial entries
the ($-$)-exponential $\Ii$-parameter map $W^-:
  \Zz_I(T,b,\mathcal{C}^-) \rightarrow W^-(T,b,\mathcal{C})$ is
  invariant under gauge transformations.
\end{prop}

\noindent{\bf Geometry and computation of $\Theta^\pm$.} The maps
$\Theta^\pm$ can be defined directly in terms of $\Ii$-parameters. A
way to see this is to consider $\Ii$ triangulations of cylinders $C =
F \times [0,1]$, as discussed in Section \ref{FIBERED}, and the
corresponding pseudo-developing maps. We give here another
description. Let us just consider $\Theta^-$. Recall that $PSL(2,\mc)$
is isomorphic to Isom$^+(\mh^3)$, with the natural conformal action on
$\mc\mathbb{P}^1=\partial\bar{\mh}^3$ via linear fractional
transformations. For any $\rho\in \Rr(g,r,\mathcal{C}^-)$, take a
representative $\tilde{\rho}: \pi_1(F,q) \rightarrow PSL(2,\mc)$ in
the conjugacy class. Consider the associated flat principal
$PSL(2,\mc)$-bundle $F_{\tilde{\rho}}$. A trivializing atlas defines a
cocycle $z \in \Zz(T,b,\mathcal{C}^-)$. If $\rho\in
\Rr_I(T,b,\mathcal{C}^-)$, we can take $z \in
\Zz_I(T,b,\mathcal{C}^-)$, so that the trivializing atlas of
$F_{\tilde{\rho}}$ associated to the cellulation $T'^*$ {\it dual} to
$T'$, with edges oriented by using the orientation of $F$ and the
branching orientation of the edges of $T'$, has non trivial transition
functions. These can be viewed as transition functions for the fiber
bundle associated to $F_{\tilde{\rho}}$ and with fibre $\mh^3$. For
each $2$-simplex $t$ of $T'$, there is a unique $g_t \in {\rm
  Isom}(\mh^3)$ (possibly reversing the orientation) mapping the
vertices $u_0$, $u_1$ and $u_2$ of the idealization of $t$ to $0$,
$\infty$ and $-1$ respectively. Then, the transition function along
the edge of $T'^*$ {\it positively transverse} to a given edge $e$ of
$T'$ is of the form $(g_{t_l})^{-1}\circ \varphi(z)(e)\circ g_{t_l}
\in {\rm Isom}^+(\mh^3)$, where $t_l$ is the triangle on the left of
$e$, and $\varphi(z)(e)$ is the isometry of $\mh^3$ of hyperbolic type
fixing $0$ and $\infty$ and mapping $1$ to $W^-(z)(e)$. Analytic
continuation defines the parallel transport of $F_{\tilde{\rho}}$
along paths transverse to $T'$, whence a representation into
$PSL(2,\mc)$ of the groupoid of such paths, well-defined up to
homotopy rel($\partial$)). In particular, it gives a practical recipe
to compute $\tilde{\rho}$, that we describe now.
\begin{figure}[ht]
\begin{center}
\includegraphics[width=10cm]{matrixrep.eps}
\caption{\label{matrixrep} The recipe for reading off holonomies
from exponential $\Ii$-parameters.}
\end{center}
\end{figure}

Let the base point $q$ be not in the $1$-skeleton of $T'$.  Given
an element of $\pi_1(F,q)$, represent it by a closed curve
$\gamma$ in $F$ transverse to $T'$, and which do not departs from
an edge it just entered. Assume that $\gamma$ intersects an edge
$e$ of $T'$ positively with respect to the orientation of $F$.
Figure \ref{matrixrep} shows three possible branching
configurations for the two triangles glued along $e$. Fix
arbitrarily a square root $W^-(z)(e)^{1/2}$ of $W^-(z)(e)$.
Consider the elements of $PSL(2,\mc)$ given by $$\gamma(e)=\left(
\begin{array}{cc} W^-(z)(e)^{1/2} & 0 \\ 0 & W^-(z)(e)^{-1/2}
  \end{array} \right), \quad p=\left(
  \begin{array}{cc} 0 & 1 \\ 1 & 0  \end{array}\right),
\quad l=\left(\begin{array}{cc} -1 & 1 \\ -1 & 0
\end{array}\right)$$
and $r=l^{-1}$. The matrix $\gamma(e)$ represents the isometry
with fixed points $0,\infty \in \mc\mathbb{P}^1$ and mapping $1$
to $W^-(z)(e)$.  The elliptic elements $p$ and $l$ send
$(0,1,\infty)$ to $(\infty,1,0)$ and $(\infty, 0,1)$ respectively.
For the portion of $\gamma$ on the left of Figure \ref{matrixrep},
if $\gamma$ turns to the left after crossing $e$ the parallel
transport operator is $\gamma(e)\cdot p\cdot l$, while it is
$\gamma(e)\cdot p\cdot r$ if $\gamma$ turns to the right.  (The
matrix multiplication is on the right, as is the action of
$PSL(2,\mc)$ on the total space of $F_{\tilde{\rho}}$). In the
middle and right pictures the parallel transport operators along
the portion of $\gamma$ are given by $\gamma(e)\cdot l$ or
$\gamma(e)\cdot p\cdot l$, and $\gamma(e)\cdot p\cdot r$ or
$\gamma(e)\cdot r$ respectively. We see that the action of $p$,
$l$ and $r$ depends on the reordering of the vertices after the
mapping $\gamma(e)$. If $\gamma$ intersects $e$ negatively, we
replace $\gamma(e)$ with $\gamma(e)^{-1}$ in the above
expressions. A similar recipe applies for any other branching of
the two triangles glued along $e$.

Continuing this way each time $\gamma$ crosses an edge of $T'$
until it comes back to $q$, we get an element $[W(z)]_q(\gamma)
\in PSL(2,\mc)$ that depends only on the homotopy class of
$\gamma$ based at $q$ and coincides with $\tilde{\rho}(\gamma)$,
because of the identity $D_z(\gamma \cdot x) = [W(z)]_q(\gamma)
\cdot D_z(x)$ for any $x \in \tilde{F}$, where $D_z: \tilde{F}
\rightarrow \mc\mathbb{P}^1$ is a $\tilde{\rho}$-equivariant
pseudo-developing map from the universal cover $\tilde{F}$ of $F$.
Varying $\gamma$, we eventually get the representation $[W(z)]_q
=\tilde{\rho}: \pi_1(F,q) \rightarrow PSL(2,\mc)$.
\smallskip

\noindent{\bf ($-$)-exponential $\Ii$-parameters and pleated
surfaces.} When $\Cc^-$ has no trivial entries, there is a nice
interpretation of the parameter space $W^-(T,b,\mathcal{C})$ in
terms of pleated hyperbolic surface structures on $S\setminus V$.
(see eg. \cite{CEG}; compare also with \cite{BoL}, Section 8.)
\smallskip

As before, let $T'$ be a triangulation of $S$ with vertices $V$,
viewed as an ideal triangulation of $S^{\circ} = S \setminus V$. A
{\it pleated surface} (with pleating locus $T'$) is a pair
$(\tilde{\mathfrak{f}},r)$, where $r: \pi_1(S^{\circ}) \rightarrow
PSL(2,\mc)$ is a group homomorphism (not up to conjugacy), and
$\tilde{\mathfrak{f}} : \widetilde{S^{\circ}} \rightarrow \mh^3$
is a map from the universal cover $\widetilde{S^{\circ}}$ of
$S^{\circ}$ such that:
\begin{itemize}
\item $\tilde{\mathfrak{f}}$ sends homeomorphically each component of
  the preimage $\tilde{T'}$ of $T'$ in $\widetilde{S^{\circ}}$ to a
  complete geodesic in $\mh^3$;

\item $\tilde{\mathfrak{f}}$ sends homeomorphically the closure of
  each component of $\tilde{S^{\circ}} \setminus \tilde{T'}$ to an
  ideal triangle in $\mh^3$;

\item $\tilde{\mathfrak{f}}$ is $r$-equivariant, that is, for all $x
  \in \tilde{S^{\circ}}$, $\gamma \in \pi_1(S^{\circ})$ we have
  $\tilde{\mathfrak{f}}(\gamma x) = r(\gamma)
  \tilde{\mathfrak{f}}(x)$.
\end{itemize}
Two pleated surfaces $(\tilde{\mathfrak{f}},r)$ and
$(\tilde{\mathfrak{f}}',r')$ are said {\it isometric} if there
exists an isometry $A \in PSL(2,\mc)$ and a lift $\tilde{\phi} :
\tilde{S^{\circ}} \rightarrow \tilde{S^{\circ}}$ of an isotopy of
$S^{\circ}$ such that $\tilde{\mathfrak{f}}' = A \circ
\tilde{\mathfrak{f}} \circ \tilde{\phi}$ and $r'(\gamma) = A
r(\gamma) A^{-1}$ for all $\gamma \in \pi_1(S^{\circ})$. From
\cite{Bo} it is known that isometry classes of pleated surfaces
$(\tilde{\mathfrak{f}},r)$ are in one-one correspondence with
arrays $\{x_e\}_e$ of non zero complex numbers $x_e$ associated to
the edges $e$ of $T'$, the {\it exponential shear-bend
parameters}.

To a pleated surface $(\tilde{\mathfrak{f}},r)$ we can associate
the type $\mathcal{C}^-=(\mathcal{C}_1,\ldots,\mathcal{C}_r)$ of
the conjugacy class of $r$, defined as in Section \ref{Dpar}. For
instance, it is shown in \cite{Th} that the set of isometry
classes of pleated surfaces given by real positive shear-bend
parameters is real-analytic diffeomorphic to the Teichm\"uller
space of hyperbolic metrics on $S^{\circ}$, with totally geodesic
boundary completion or non compact finite area completion at
$v_i$, according to the type of $\mathcal{C}_i$, loxodromic or
parabolic.

In general, we can also associate a {\it sign} to each loxodromic
puncture. Namely, if $A$ is a small annulus neignborhood of $v_i$,
each connected component $\tilde{A}$ of the preimage of $A$ in
$\tilde{S^{\circ}}$ is the fixed point set of a subgroup $\pi$ of
$\pi_1(S^{\circ})$. This subgroup is the image of the fundamental
group of $A$ for some choice of base points and paths between
these base points. All the edges of $\tilde{T}'$ that meet
$\tilde{A}$ are sent by $\tilde{\mathfrak{f}}$ to geodesics lines
that meet at one of the two fixed points of $r(\pi)$. We specify
this fixed point by a sign, as it determines an orientation
(whence a generator) for the axis of the group $r(\pi)$. Since any
two subgroups $\pi$ as above are conjugated, for each puncture the
fixed point assignment is $r$-equivariant, so that the sign is
canonically associated to the puncture.

\smallskip

Recall that in Definition \ref{defcoc} the types $\mathcal{C}_i$
where assigned to the components of $\partial F$ endowed with the
boundary orientation. Let us remove this constraint, and associate
to each boundary component with loxodromic type an orientation
that we specify by a sign, positive for the boundary orientation,
and negative otherwise. Then, for each type $\mathcal{C}^-$ with
$l$ loxodromic entries and each $l$-uple of signs $s$, we get a
space $\Zz_I(T,b,\mathcal{C}^-,s)$. We define
$W^-(T,b,\mathcal{C},s)=W(\Zz_I(T,b,\mathcal{C}^-,s))$.
\begin{prop}\label{psI}{\rm For each type
    $\mathcal{C}$ with non trivial entries $\mathcal{C}_i$, $l$ being
    loxodromic, and for each $l$-uple of signs $s$, the space
    $W^-(T,b,\mathcal{C},s)$ coincides with the exponential shear-bend
    parameter space $\mathcal{P}\mathcal{S}_{\mathcal{C}}^s$ of
    isometry classes of pleated surfaces of type $\mathcal{C}$ and
    signs $s$:
\begin{equation}\label{polPS}
\mathcal{P}\mathcal{S}_{\mathcal{C}}^s = \lbrace \{x_e\} \in (\mc
  \setminus \{ 0 \})^{-3\chi(S^{\circ})} \mid \forall i=1,\ldots,r,\
  \prod x_e = \mu_i \rbrace.
\end{equation}
Here the product is over all edges $e$ with $v_i \in\partial e$
(counted with multiplicities), $\mu_i$ is $1$ if $\mathcal{C}_i$
is parabolic, and, if $\mathcal{C}_i$ is loxodromic, $\mu_i$ is
the dilation factor of the generator of $<\mathcal{C}_i>$
specified by the sign of $v_i$. (Hence, in either case this is an
eigenvalue of $\mathcal{C}_i$.)}
\end{prop}

\noindent{\it Proof.} By the results of \cite{Bo} recalled above,
each point of $W^-(T,b,\mathcal{C},s)$ is identified with the
family of exponential shear-bend parameters of an isometry class
of pleated hyperbolic surface on $S^{\circ}$ with pleated locus
$T'$. The definition of the map $\Theta: W^-(T,b,\mathcal{C},s)
\rightarrow \Rr_I(T,b,\mathcal{C})$ shows that this isometry class
has type $\mathcal{C}$. Furthermore, by using the recipe given
above for computing $\Theta$, we check that the holonomy of any
positively oriented (with respect to the boundary orientation)
small loop about the puncture $v_i$ is exactly $\mathcal{C}_i$.
The upper left diagonal entry is just the product of a square root
of the exponential $\Ii$-parameters at the edges with endpoint
$v_i$. From (\ref{polPS}), which is an easy consequence of results
in \cite{Bo} (see sections 12.2-12.3 in that paper), we deduce
that $W^-(T,b,\mathcal{C},s) \subset
\mathcal{P}\mathcal{S}_{\mathcal{C}}^s$.

Conversely, for any $(\tilde{\mathfrak{f}},r) \in
\mathcal{P}\mathcal{S}_{\mathcal{C}}^s$ we have $ {\rm conj}(r) \in
\Rr_I(T,b,\mathcal{C})$, because $r$ is injective. Also, by
\cite{BoL}, Proposition 33, it is known that the isometry class of
$(\tilde{\mathfrak{f}},r)$ is determined by $r$ and the signs $s$.
As the map $\Theta$ is onto, let us take $z \in
\Zz_I(T,b,\mathcal{C}^-,s)$ with ${\rm conj} \circ f(z) = {\rm conj}(r)$
and consider the isometry class of pleated hyperbolic surfaces
associated to $W^-(z)$. There is a representative
$(\tilde{\mathfrak{f}}',r)$ with the same holonomy $r$. As the signs
of $(\tilde{\mathfrak{f}}',r)$ are $s$, the same as for
$(\tilde{\mathfrak{f}},r)$, the two pleated surfaces coincide. So
$\mathcal{P}\mathcal{S}_{\mathcal{C}}^s \subset
W^-(T,b,\mathcal{C},s)$.\hfill$\Box$

\begin{remark}\label{branchwhy}{\rm
Exponential shear-bend parameters do not depend on branchings but the
orientation of $F$ (taking the opposite branching edge orientation
simultaneously exchanges the left and right triangles). However,
branchings govern all choices in QHFT tensors. Also, they allow us to
interpret exponential shear-bend parameters as exponential
$\Ii$-parameters, thus coming from $1$-cocycles representing {\it
  arbitrary} $PSL(2,\mc)$-characters on the triangulated boundary of
{\it arbitrary} compact orientable $3$-manifolds.  When $\Cc$ has no
trivial entries, the maps $W^\pm$ define {\it decorated} shear-bend
parameter spaces similar to those occurring in \cite{Kmod, Pe}.}
\end{remark}
\begin{remark}\label{varpar}
{\rm For types $\mathcal{C}$ with no trivial entries, we can use
simpler $e$-triangulations of $F$ (see Figure \ref{more-e}): for
each of the triangles $t_v$ of the base surface $S$ we remove the
interior of a {\it monogon} inside $t_v$, and triangulate the
resulting quadrilateral by adding an edge $e_v$ with endpoints $v$
and the $b$-output vertex of the opposite edge. We extend $b$ by
orienting $e_v$ from that vertex to $v$. Proposition \ref{topcomp}
applies to the cocycle parameters based on such
$e$-triangulations, for points of $\Rr(g,r)$, and also the
treatment of exponential $\Ii$-parameters works as well.}
\end{remark}
\begin{figure}[ht]
\begin{center}
\includegraphics[width=3.5cm]{moreE.eps}
\caption{\label{more-e} More economic e-triangulation.}
\end{center}
\end{figure}
